% Kenny & Pearce: Brexit and the Anglosphere: Reading Summary
% Introduction to the European Union, Week 5
% Compile with: pdflatex kenny-pearce-brexit-anglosphere-reading-summary.tex

\documentclass[11pt,a4paper]{article}
\usepackage[utf8]{inputenc}
\usepackage[T1]{fontenc}
\usepackage[margin=2.5cm]{geometry}
\usepackage{booktabs}
\usepackage{enumitem}
\setlist{nosep,leftmargin=*}
\usepackage{parskip}
\usepackage{microtype}
\usepackage{hyperref}
\usepackage{xcolor}
\definecolor{eublue}{RGB}{0,51,153}

\hyphenation{Spitzenkandidat Spitzenkandidaten}
\hypersetup{
  colorlinks=true,
  linkcolor=eublue,
  urlcolor=eublue,
  pdftitle={Kenny \& Pearce: Brexit and the Anglosphere},
  pdfauthor={Introduction to the EU, UCM}
}

\begin{document}

\title{Kenny \& Pearce: Brexit and the Anglosphere}
\author{Introduction to the European Union\\Universidad Complutense de Madrid\\Week 5\\In Stuart Ward and Astrid Rasch (eds.), \textit{Embers of Empire in Brexit Britain}, Bloomsbury 2019}
\date{}
\maketitle
\vspace{-1em}
\hrule
\vspace{1em}

\section*{Rhetorical Framework (Ethos, Logos, Pathos)}

\begin{table}[h]
\centering
\begin{tabular}{@{}lll@{}}
\toprule
\textbf{Appeal} & \textbf{Meaning} & \textbf{Kenny \& Pearce's Use} \\
\midrule
\textbf{Ethos} & Credibility, authority, trust & Historical depth; moves beyond short-term explanations; scholarly synthesis \\
\textbf{Logos} & Logic, structure, argument & Anglosphere vs.\ Europe as identity tension; imagined community as analytical category \\
\textbf{Pathos} & Emotion, stakes, persuasion & Empire, memory, sovereignty; emotional appeal of Anglosphere narrative \\
\bottomrule
\end{tabular}
\end{table}

\textbf{Key insight:} Kenny \& Pearce show how pro-Brexit actors deploy \textit{pathos} (empire, Anglosphere, sovereignty) to build \textit{ethos} (British identity, natural belonging) and challenge \textit{logos} of EU membership. The Anglosphere is ``imagined'' but \textit{rhetorically} powerful.

\section*{Bibliographic Information}

\begin{itemize}
\item \textbf{Authors:} Mike Kenny and Nick Pearce
\item \textbf{Title:} Brexit and the Anglosphere
\item \textbf{Source:} In Stuart Ward and Astrid Rasch (eds.), \textit{Embers of Empire in Brexit Britain}, Bloomsbury Publishing
\item \textbf{Year:} 2019
\item \textbf{Pages:} 25--35
\end{itemize}

\section*{Summary \& Key Points}

\begin{itemize}
\item The authors connect Brexit to long-standing British ideas about belonging to an ``Anglosphere''---a community of English-speaking nations (UK, US, Canada, Australia, etc.) anchored in imperial history.
\item They show how pro-Brexit actors mobilised this Anglosphere narrative as an attractive alternative to European integration.
\item The chapter emphasises how memories of empire and wartime victory shape contemporary British understandings of sovereignty and international status.
\end{itemize}

\section*{Main Arguments}

\subsection*{I. British Identity: Europe vs.\ Anglosphere}

Kenny and Pearce argue that Brexit is partly the result of a deep tension in British identity: between being a European state and being the centre of a wider Anglosphere. Brexit revived the latter imagination, presenting Europe as a constraint and the English-speaking world as a natural sphere of influence.

\subsection*{II. Anglosphere as Imagined Community}

They also contend that the Anglosphere is as much an imagined community as a concrete geopolitical project. It is powerful in rhetoric and symbolism, even if its economic and political foundations are uncertain.

\section*{Context \& Analysis}

``Anglosphere'' = \textit{pathos} (belonging, identity) + \textit{ethos} (British exceptionalism). ``Imagined community'' = \textit{logos} (analytical category). Pro-Brexit rhetoric as case study in how narratives shape politics.

\section*{Relevance to Course}

\begin{itemize}
\item \textbf{Ethos:} Explains why the UK has often seen itself as a ``reluctant European'' even before 2016.
\item \textbf{Logos:} Anglosphere vs.\ Europe as analytical framework; imagined community as concept.
\item \textbf{Pathos:} Euroskepticism entangled with imperial memory and global identity, not only EU policies.
\item \textbf{All three:} Different historical narratives compete when defining a country's place in or out of the EU.
\end{itemize}

\section*{Discussion Questions}

\begin{itemize}
\item How does Anglosphere narrative challenge European project? (Ethos: competing identities; Logos: structural tension; Pathos: emotional appeal of each)
\item To what extent is Brexit a post-imperial adjustment crisis? (Ethos: who frames it? Logos: evidence, causes; Pathos: stakes of identity)
\item Can the EU offer an alternative narrative to compete with Anglosphere? (Ethos: who tells the EU story? Logos: what narrative resources? Pathos: emotional appeal, belonging)
\end{itemize}

\vspace{1em}
\noindent\footnotesize\textit{Reference:} Kenny, Mike and Nick Pearce (2019). ``Brexit and the Anglosphere.'' In Stuart Ward and Astrid Rasch (eds.), \textit{Embers of Empire in Brexit Britain}, London: Bloomsbury, pp. 25--35.

\end{document}
